\documentclass[11pt]{article}
\usepackage[T1]{fontenc}
\usepackage{lmodern}
\usepackage[margin=0.95in]{geometry}
\usepackage{amsmath,amssymb,amsthm,booktabs,microtype}
\usepackage[colorlinks=true,linkcolor=blue,citecolor=blue,urlcolor=blue]{hyperref}
\hypersetup{pdftitle={Preparation dependence of proper-time volume in Einstein-scalar shell perturbations},pdfauthor={Ricardo Maldonado},pdfsubject={Working manuscript for author and specialist review}}
\pdfinfoomitdate=1
\pdftrailerid{}
\pdfsuppressptexinfo=15
\newcommand{\eps}{\varepsilon}
\newcommand{\dd}{\mathrm d}
\newcommand{\Lie}{\mathcal L}
\newcommand{\av}[1]{\langle #1\rangle}
\newcommand{\norm}[1]{\lVert #1\rVert}
\newcommand{\cV}{\mathcal V}
\newtheorem{theorem}{Theorem}
\newtheorem{lemma}{Lemma}
\title{Preparation dependence of proper-time volume in Einstein--scalar shell perturbations}
\author{Ricardo Maldonado\\\small Independent Researcher}
\date{4 November 2026}

\begin{document}
\maketitle
\begin{center}\small Working manuscript for author and specialist review\end{center}
\begin{abstract}
A first perturbation and its quadratic sources need not determine the sign
of a second observable when physical homogeneous second initial data remain
free. We give a conditional quantitative example for a five-dimensional
Einstein--scalar spacetime with a scalar-dependent pure-tension shell.
A compact zero-moment metric mean changes the second preparation while
preserving the first quadrupole packet, canonical constraints, initial shell
collar, and second transverse-traceless initial pair. The difference solves
the complete homogeneous five-mean problem. At one fixed anchored proper
time, its proper-volume and expansion slopes are strictly negative.
The fixed finite preparation parameter $\lambda=-10^{21}$ therefore changes
both second coefficients from negative to positive. Whole-cell rational
enclosures give $v_2\in[1.0343,1.8442]\times10^{17}$ and
$\Theta_2\in[8.3831,18.496]\times10^{21}$; a separate elementary bound
also proves both signs. The result assumes the stated constraint-inverse,
development, regularity, and linear-uniqueness premises. It supplies no
effective nonlinear parameter radius, contracting branch, or physical rule
selecting a preparation. The affine structure is established perturbation
machinery; the contribution is this model-specific signed example.
\end{abstract}

\section{Model, fixed data, and conditional result}

Bulk initial information can survive in brane observables
\cite{Minamitsuji,Deffayet,Hiramatsu}. Ordinary second-variation theory
also separates quadratic forcing from homogeneous second data. The question
here is more specific: can two compact, constraint-compatible preparations
with the same first packet and initial shell collar give opposite signs
for both proper-volume and expansion second coefficients at the same event?
Non-gauge data alone do not answer it; their observable response must be
shown to be nonzero. We construct one direction and bound that response.

The bulk and one-copy outward shell equations are
\begin{align}
 \operatorname{Ric}_{AB}-\phi_A\phi_B-\tfrac23\cV(\phi)g_{AB}&=0,
 &\Box_g\phi-\cV'(\phi)&=0, \label{eq:model}\\
 K^{\rm out}_{ab}&=\tfrac16\sigma(\phi)\gamma_{ab},
 &n_{\rm out}\phi&=-\tfrac12\sigma'(\phi). \label{eq:shell}
\end{align}
Two identical interior copies use their respective outward normals. Fix
\begin{gather}
 W(\phi)=1-\phi+\phi^3/3,\qquad
 \cV=\tfrac12(W')^2-\tfrac23W^2,\qquad
 \sigma=2W+\delta(1+c_\sigma\phi),\nonumber\\
 \delta=10^{-3},\qquad c_\sigma=96714321425692/161839258930731.
 \label{eq:potential}
\end{gather}
The selected near-$\phi=1$ regular background is
\begin{equation}
 g_0=\dd y^2+\rho(y)^2[-\dd t^2+\cosh^2t\,\kappa],\qquad
 a=\rho_y/\rho,\quad p=\phi_{0,y}, \label{eq:background}
\end{equation}
where $\kappa$ is the unit round metric on $S^3$. Its exact selected shell
radius is $\rho_b$; it is not a rounded value chosen from a bounding box.
Put $\ell=10^{-6}$, $x=t/\ell$, $r=z/\ell$, and
$\dd y=-\rho\,\dd z$. The packet support is $1\le r\le2$.
The fixed target is $T=5/4$, or $\tau_*=\rho_b\ell T$ in background
proper time. All quantities below are in these model units.

We use second parameter derivatives throughout:
\begin{equation}
 x_\eps=x_0+\eps v_{\rm seed}+\tfrac12\eps^2x_2+o(\eps^2),\qquad
 g_\eps=g_0+\eps q+\tfrac12\eps^2k+o(\eps^2).
 \label{eq:factorial}
\end{equation}
Anchor the shell's intrinsic future normal-geodesic foliation at the
common initial section, with induced spatial metric $h_\eps(\tau)$, and set
\begin{equation}
 V_\eps=\int_{S^3}\dd\mu_{h_\eps},\quad
 \Theta_\eps=\partial_\tau\log V_\eps,\quad
 M_\eps=V_\eps^{-1}\int_{S^3}\phi_\eps\,\dd\mu_{h_\eps},\quad
 v_2=V_2/V_0. \label{eq:observables}
\end{equation}
The clock is intrinsic to the shell and transported under gauge changes.
The background has $A(\tau)=\rho_b\cosh(\tau/\rho_b)$,
$H=A_\tau/A$, and $V_0=2\pi^2A^3$.

The analytic premises are explicit. We assume the selected regular
background and its authenticated annulus bounds; the equivariant weighted
constraint isomorphism and nonlinear implicit-function hypotheses;
the common fixed-shell development and subsidiary/boundary uniqueness;
and identification of its first derivative with the retained first-field
solution. For each fixed finite preparation, assume $C^2$ parameter
dependence in the retained lower $C^4$ spacetime topology, with the traces
and derivatives required in \eqref{eq:observables}. A common intrinsic
normal tube reaches the target after a family-dependent qualitative
shrinkage in $\eps$. The full linear mean mixed problem is unique and
obeys the coefficient bounds in Section~\ref{sec:bounds}. These analytic
premises are retained from the technical development \cite{records}; they
are not proved by finite interval arithmetic. No effective admissible
$\eps$ radius is assumed or obtained.

\begin{theorem}[Conditional preparation dependence]\label{thm:main}
Under these premises, there are two exact completed families with the
same background, first packet, canonical quadratic constraint source,
initial shell collar, and second transverse-traceless initial pair, for
which the second coefficients at the same proper clock satisfy
\begin{center}
\begin{tabular}{@{}lcc@{}}\toprule
Preparation & $v_2(\tau_*)$ & $\Theta_2(\tau_*)$\\\midrule
Base & $[-5.4520,-2.6521]\,10^{16}$ & $[-9.0208,-4.4458]\,10^{21}$\\
$\lambda=-10^{21}$ & $[1.0343,1.8442]\,10^{17}$ & $[8.3831,18.496]\,10^{21}$\\
\bottomrule
\end{tabular}
\end{center}
These are outward decimal displays of exact rational enclosures.
Consequently the retained first-packet specification does not determine
either sign. This single direction proves simultaneous reversal, without
asserting independent control of the two coefficients.
\end{theorem}

\section{A compact constraint-compatible preparation direction}

Here are the fixed first data and the base preparation, so that the
comparison identifies what is held constant. In a background orthonormal
spatial frame, with background spatial density $\rho_0$, the first
canonical data for the action normalization
$R/2-|\dd\phi|^2/2-\cV$ are
\begin{equation}
 \gamma_1=\operatorname{diag}(-4u,2u,2u,2u)Y,\quad
 \pi_1/\rho_0=\operatorname{diag}(-3s_1/2,0,0,0)Y,\quad
 \phi_1=vY,\quad\varpi_1/\rho_0=jY. \label{eq:first}
\end{equation}
The angular factor is $Y=n_1^2-n_2^2$, with
$\av Y=0$, $\av{Y^2}=1/6$, and $\Delta_\kappa Y=-8Y$; brackets use
probability measure. Define
\begin{align}
 Q(r)&=\exp\!\left(4-\frac1{(r-1)(2-r)}\right)\quad(1<r<2),
 \qquad Q=0\text{ otherwise},\label{eq:Q}\\
 B(z)&=Q(z/\ell),\quad
 \alpha=(\rho/\rho_b)^{3/2}p/p_b,\quad
 \varkappa=\alpha_z/\alpha=\rho(5a/2-\cV'/p),\nonumber\\
 u&=\alpha B/\rho^2,\quad u_t=e^{-z}u,\quad s_1=u_t/\rho,\nonumber\\
 v&=3\alpha(B_z+\varkappa B)/(\rho^3p),\quad
 v_t=3\alpha e^{-z}[B_z+(\varkappa-1)B]/(\rho^3p),\quad j=v_t/\rho.
 \nonumber
\end{align}
All endpoint extensions are smooth and flat. In particular
$pv=-3(u_y+2au)$ and $pj=-3((s_1)_y+3as_1)$.
The mean Hamiltonian and radial momentum Hessians are
\begin{align}
 Q_H={}&-2u_y^2-24auu_y-2uu_{yy}
 -(16a_y+32a^2+8p^2/3)u^2-4u^2/\rho^2-s_1^2\nonumber\\
 &-(j^2+v_y^2+\cV''v^2)/6-(4/3)puv_y-(4/3)v^2/\rho^2,
 \label{eq:QH}\\
 Q_M={}&-2s_1u_y+j(v_y+4pu)/3.\nonumber
\end{align}
They are Hessians in convention \eqref{eq:factorial}. Let
$y_L=y(2\ell)<y_U=y(\ell)$ and
\begin{align}
 I&=\int_{y_L}^{y_U}\rho^4Q_H\,\dd y,\quad
 J=\int_{y_L}^{y_U}\rho^2 B\,\dd y>0,\quad f=\frac{2I}{9J}B,
 \nonumber\\
 C[f]&=-\tfrac32(f_{yy}+3af_y)-\cV f,\qquad
 \xi=-\frac{\int_{y_L}^y\rho^4(Q_H+C[f])\,\dd s}{\rho^4p^2}.
 \label{eq:base}
\end{align}
The base second mean is the canonical datum
\begin{equation}
 B_{\rm base}=\left(f\gamma_0+\Lie_{\xi\partial_y}\gamma_0,
 \rho_0\operatorname{diag}(0,-Q_M/(6a),-Q_M/(6a),-Q_M/(6a)),0,0\right).
 \label{eq:Bbase}
\end{equation}
It solves the second mean constraint equation $PB_{\rm base}=-P_0\mathcal Q$,
where $\mathcal Q$ is the complete first constraint Hessian and $P_0$ is
Haar averaging. The base numerical response in Theorem~\ref{thm:main}
retains the actual first evolution and shell forcing, not merely
\eqref{eq:QH}.

Now let $h$ be compact and flat in the same packet interval, with
\begin{equation}
 \int_{y_L}^{y_U}\rho^2h\,\dd y=0,\qquad
 h=\rho^{-2}\chi_y,\qquad
 \chi(y)=\int_{y_L}^y\rho^2h\,\dd s. \label{eq:moment}
\end{equation}
For every fixed finite real $\lambda$, set
\begin{equation}
 f_\lambda=f+\lambda h,\quad
 \xi_\lambda=\xi+\lambda\zeta,\qquad
 \zeta=\frac{3(h_y-ah)}{2p^2}+\frac{9\chi}{2\rho^4p^2}.
 \label{eq:zeta}
\end{equation}
The canonical momentum mean and both scalar canonical sectors in
\eqref{eq:Bbase} are unchanged. Thus
$B_\lambda-B_{\rm base}=\lambda(H_h,0,0,0)$, where
$H_h=h\gamma_0+\Lie_{\zeta\partial_y}\gamma_0$.
The background identity
$a_y+4a^2=3/\rho^2-2\cV/3$ gives the exact flux formula
\begin{equation}
 \rho^4C[h]=-\tfrac32[\rho^4(h_y-ah)]_y-\tfrac92\rho^2h,
 \quad
 P_gH_h=C[h]+\rho^{-4}(\rho^4p^2\zeta)_y=0. \label{eq:flux}
\end{equation}
The moment condition makes its primitive zero on both exterior sides.
Flatness and $p\ne0$ give compact smooth extension. This changes no part
of the quadratic source: $PB_\lambda=-P_0\mathcal Q$ exactly.

For fixed retained weights, insert $\eps^2B_\lambda/2$ into the same
weighted implicit-function trial
$x_\eps=x_0+\eps v_{\rm seed}+\eps^2B_\lambda/2+\mathcal A U_\lambda(\eps)$.
The correction derivative is the same isomorphism
$L=\omega^{-2}P\mathcal A$, and
\begin{equation}
 LU_\lambda''(0)=-\omega^{-2}(\mathcal Q+PB_\lambda).
 \label{eq:ift}
\end{equation}
It is independent of $\lambda$. Retained equivariance and injectivity
give $P_0U_\lambda''(0)=0$, hence exactly the declared second mean.
Compact support preserves the initial shell collar and second TT pair.
The common first tangent and retained linear uniqueness preserve the first
field and all its quadratic sources. For every fixed finite $\lambda$,
the existence and support arguments apply after qualitative shrinkage in
$\eps$. No radius uniform as $|\lambda|$ grows, no choice
$\lambda=\lambda(\eps)$ tending to infinity, and no affine full
finite-$\eps$ family are claimed.

The freedom is physical in the anchored pure second gauge quotient with
the common first representative fixed. On the time-symmetric initial
section $K_0=\Pi_0=0$, so a normal generator changes neither initial
spatial metric nor scalar. Averaging a proposed tangential generator leaves
a radial generator; there is no invariant angular one-form. Its scalar
variation is $pX^y$. The raw scalar difference is zero, so $X^y=0$.
For completeness, $H_h=0$ would imply
$h+2\zeta_y=h+2a\zeta=0$, hence $\zeta_y=a\zeta$; compact support then
gives $\zeta=h=0$. A nonzero $h$ is therefore not such a diffeomorphism.
This does not imply that every observable detects it.

\section{Conditioning, proper observables, and exact affine response}

For each preparation a pure second harmonic vector has initial data
$(Z^y,Z^t,Z^y_t,Z^t_t)=(-\xi_\lambda,0,0,a\xi_\lambda)$ and shell
rows $Z^y=0$, $\partial_yZ^t=0$. It solves
$(\Box_0+\operatorname{Ric}_0)Z=0$ and transports the shell anchor and
clock. The transformation
$\bar k=k+\Lie_Zg_0$, $\bar w=w+Z\phi_0$ removes the spatial Lie term
and keeps its essential scalar pullback. It is a change of representative
of each family, not the physical preparation freedom just constructed.
The second jet suffices; no nonlinear harmonic transformation at finite
$\eps$ is asserted. The factorial convention and transport follow the
standard second-order gauge framework \cite{Bruni,Mars}.

In the background orthonormal frame, let $(K,L,N,J,w)$ denote the
$yy,y0,00$, angular-isotropic, and scalar mean slots. Subtracting the
conditioned adapters gives the initial data for the unit direction:
\begin{equation}
 \begin{aligned}
 \mathbf W(0,r)&=(h,0,0,h,s),\qquad s=-p\zeta,\\
 \mathbf W_x(0,r)&=(0,\ell\rho(-h_y+ah),0,0,0),\qquad
 \mathbf W_r(0,r)=(h_r,0,0,h_r,s_r).
 \end{aligned}\label{eq:initialW}
\end{equation}
The mixed metric velocity is retained. The scalar coordinate velocity is
zero; the first-only lapse-square term cancels in this difference.
The actual forced mean system has identical operator and shell rows for
both preparations. Uniqueness therefore gives
\begin{equation}
 \mathbf U_\lambda=\mathbf U_{\rm base}+\lambda\mathbf W,\qquad
 \mathbf W_{xx}-\mathbf W_{rr}=A_0\mathbf W_x+B_0\mathbf W_r+D_0\mathbf W,
 \quad L|_{r=0}=0,\quad P_N\mathbf W_r=M_0P_N\mathbf W. \label{eq:W}
\end{equation}
Here $P_N$ selects the other four slots. All metric--scalar, connection,
physical Robin, and normal-gauge couplings are included.

The proper scalar mean illustrates why absolute and difference responses
must be distinguished. In the background orthonormal shell frame
$e_0=\partial_\tau$, $e_a=A^{-1}E_a$, with $E_a$ unit-round, write the
first shell metric components as
$m_{00}=-2XY$, $m_{0a}=\mathsf V D_aY$,
$m_{ab}=2(ZY\kappa_{ab}+dD_aD_bY)$ and the first scalar as $\psi Y$.
With $C_\tau=X$, $C(0)=0$, and $E=Z-HC$, its second proper mean is
\begin{equation}
 M_2=w_0-C\psi_\tau/3+(E-8d/3)\psi. \label{eq:propermean}
\end{equation}
The last two terms are clock and volume-measure corrections. Every first
quantity here is shared; these terms cancel only on subtraction.

The shell Gauss contraction \cite{SMS} gives
$R(\gamma)=(D\phi)^2+F(\phi)$, where
$F=2\cV-(\sigma')^2/4+\sigma^2/3$. In the anchored foliation its
volume-weighted second identity is
\begin{align}
 (v_2)_{\tau\tau}+4H(v_2)_\tau-2v_2/A^2
   &=\tfrac12F'_bM_2+Q_1,\qquad v_2(0)=(v_2)_\tau(0)=0,
 \label{eq:volume}\\
 Q_1&=F''_b\psi^2/12-\psi_\tau^2/6+4\psi^2/(3A^2)-5E^2/(3A^2)
 \nonumber\\
 &\hspace{12mm}+E_\tau^2-16E_\tau D_\tau/3+8D_\tau^2/3,
 \qquad D_\tau=d_\tau-\mathsf V/A-C/A^2.\nonumber
\end{align}
Thus every first-only quadratic term $Q_1$ also cancels in the difference.
The absolute base response retains the full forcing, the corrections in
\eqref{eq:propermean}, and every term of $Q_1$. No source is removed
from that response to obtain the positive witness.

Put $m(x)=W_w(x,0)$. Equations \eqref{eq:propermean}--\eqref{eq:volume}
give the exact affine identities
\begin{equation}
 M_{2,\lambda}=M_{2,\rm base}+\lambda m,\qquad
 v_{2,\lambda}=v_{2,\rm base}+\lambda\nu,\qquad
 \Theta_{2,\lambda}=\Theta_{2,\rm base}+\lambda\theta,\quad
 \theta=\nu_\tau. \label{eq:affine}
\end{equation}
There is no extra factor $1/2$ in these second-derivative slopes.
The exact target functionals are
\begin{align}
 \mathcal L_M[h]&=m(T),\nonumber\\
 \mathcal L_v[h]&=(\rho_b\ell)^2(F'_b/2)\int_1^T k(T,x)m(x)\,\dd x,
 \label{eq:slopes}\\
 \mathcal L_\Theta[h]&=\rho_b\ell(F'_b/2)\int_1^T k_T(T,x)m(x)\,\dd x,
 \nonumber\\
 k(T,x)&=\frac{\sinh(2\ell(T-x))}{2\ell}
              \frac{\cosh^2(\ell x)}{\cosh^2(\ell T)}.\nonumber
\end{align}
$F'_b<0$ on the selected background. The homogeneous solutions
$\tanh(\ell x)$ and $1+\tanh^2(\ell x)$ give this kernel. For
$v=\ell T$, $b=\ell x$,
$k_T=\cosh^2b\cosh(v-2b)/\cosh^3v$. Therefore
\begin{equation}
 0\le k\le T-x,\quad 0<k_T\le1,\qquad
 1-\eta\le k/(T-x)\le1,\quad 1-\eta\le k_T\le1,
 \quad \eta=10^{-11}. \label{eq:kernelcaps}
\end{equation}
The ratio extends continuously at $x=T$. Lower bounds
$(1+v^2)^{-2}$ and $(1+v^2)^{-3}$ suffice since $3v^2<\eta$.

\section{A full coupled response certificate}\label{sec:bounds}

Fix the single direction
\begin{equation}
 \chi=\ell\rho_b^3Q(r),\qquad h=-(\rho_b/\rho)^3Q_r(r).
 \label{eq:direction}
\end{equation}
It satisfies \eqref{eq:moment} identically, without numerical moment
cancellation. Define
\begin{equation}
 R=(\rho_b/\rho)^3,\quad b=\ell\rho a,\quad
 c=-\frac{3\rho_b^3}{2\ell\rho^4p}>0,\quad d_0=\ell\rho p_y/p.
 \label{eq:profileconstants}
\end{equation}
Primes in the following formulas mean $r$ derivatives:
\begin{align}
 h&=-RQ',\quad h_r=-R(Q''+3bQ'),\quad
 V_L=-R(Q''+4bQ'),\nonumber\\
 s&=c(Q''+4bQ'+3\ell^2Q),\label{eq:profiles}\\
 s_r&=c\{Q'''+4bQ''+3\ell^2Q'
 +(4b+d_0)(Q''+4bQ'+3\ell^2Q)
 -4\ell^2\rho^2(a^2+a_y)Q'\}.\nonumber
\end{align}
These use $\rho_r=-\ell a\rho^2$, $R_r=3bR$,
$c_r=c(4b+d_0)$, and $b_r=-\ell^2\rho^2(a^2+a_y)$.
The initial Gram norms for $\mathsf H=\operatorname{diag}(1,2,1,3,1)$ are
$|\mathbf W_0|_{\mathsf H}=\sqrt{4h^2+s^2}$,
$|\mathbf W_{0,r}|_{\mathsf H}=\sqrt{4h_r^2+s_r^2}$, and
$|\mathbf W_{0,x}|_{\mathsf H}=\sqrt2|V_L|$.

On $D_T=\{0\le x\le T,\ 0\le r\le T-x\}$ use
\begin{equation}
 X(\mathbf U)=\max\{\sup|\mathbf U|_{\mathsf H}/3,
                 \sup|\mathbf U_x|_{\mathsf H},\sup|\mathbf U_r|_{\mathsf H}\}.
 \label{eq:X}
\end{equation}
The retained coefficient caps are
$\norm{A_0},\norm{B_0}\le10^{-4}$, $\norm{D_0}\le10^{-6}$,
$\norm{M_0}\le1/2000$. With $\Delta=T-1=1/4$, the reflected bulk
Duhamel operator and Robin lift have bounds $T\sup|F|_{\mathsf H}$
and $\sup|G|_{\mathsf H}$. Their combined contraction is
\begin{equation}
 q=T(2\cdot10^{-4}+\Delta10^{-6})+\Delta/2000
   =1201/3200000<1. \label{eq:q}
\end{equation}
Let $\mathbf B$ be the free reflected wave of \eqref{eq:initialW} and
let $D_0^{\rm init},D_r^{\rm init},D_v^{\rm init}$ bound its initial
value, radial derivative and velocity Gram norms on $0\le r\le T$.
Then
\begin{equation}
 N=\max\{(D_0^{\rm init}+TD_v^{\rm init})/3,
                  D_r^{\rm init}+D_v^{\rm init}\}\ge X(\mathbf B),\qquad
 E_h=\frac{qN}{1-q}\ge X(\mathbf W-\mathbf B). \label{eq:error}
\end{equation}
This is the full resolvent, not a truncated iteration or an uncoupled
scalar approximation. The common actual forcing cancels exactly in
\eqref{eq:W}; no source-history budget is needed for this difference.

Both waves vanish in $x+r\le1$. At the shell, even scalar reflection and
zero scalar initial velocity give $B_w(x,0)=s(x)$. Hence
\begin{equation}
 |m(x)-s(x)|\le(x-1)_+E_h,\qquad
 |m(T)-s(T)|\le\Delta E_h. \label{eq:shellerror}
\end{equation}
Define free moments $J_v=\int_1^T(T-x)s(x)\,\dd x$ and
$J_\Theta=\int_1^Ts(x)\,\dd x$, with absolute versions $A_v,A_\Theta$.
The kernel-weighted histories differ from these by at most
\begin{equation}
 \frac{\Delta^3E_h}{6}+\eta A_v,\qquad
 \frac{\Delta^2E_h}{2}+\eta A_\Theta, \label{eq:historyerror}
\end{equation}
respectively. This proves continuous-history bounds without assuming a
pointwise sign of the evolved scalar arbitrarily close to first contact.

The registered interval calculation uses the fixed $N=1024$ partition
restricted to this causal initial segment: one grouped flat tail
$[1,65/64]$ and 240 interior cells
$[k/1024,(k+1)/1024]$, $k=1040,\ldots,1279$. It makes 241 whole-cell
bump evaluations and one distinct bump target evaluation at $r=5/4$,
242 in total, with zero new background or evolved-first-field oracle
calls. The segment $[0,1]$ is exactly zero. The flat tail uses the
monotone bounds for $s_0^{2j}e^{4-s_0}$ with
$s_0\ge4096/63$, $j\le3$, including the flat endpoint. Positive exact
cell weights integrate the whole-cell scalar boxes. Fresh profile bounds
are used; old preparation profiles and old total-response norms are not
substituted for \eqref{eq:profiles}.

Outward upper displays are
\begin{equation}
 D_0^{\rm init}\le3.0797\,10^9,\quad D_r^{\rm init}\le4.5253\,10^{10},
 \quad D_v^{\rm init}\le46.892,\quad N\le4.5253\,10^{10},
 \quad E_h\le1.6991\,10^7. \label{eq:numericnorms}
\end{equation}
The exact rational arithmetic in \eqref{eq:slopes}--\eqref{eq:historyerror}
gives
\begin{center}
\begin{tabular}{@{}lc@{}}\toprule
Full coupled slope & Outward enclosure\\\midrule
$\mathcal L_M[h]$ & $[1.5393,1.7068]\,10^9$\\
$\mathcal L_v[h]$ & $[-2.1094,-1.5795]\,10^{-4}$\\
$\mathcal L_\Theta[h]$ & $[-22.9410,-17.4038]$\\\bottomrule
\end{tabular}
\end{center}
The nonzero observable slopes are thus certified, rather than inferred
from non-gauge initial data. Normal and optimized runs and an independent
reconstruction agree in their rational scientific outputs \cite{records}.
They share analytic premises and are not independent observations of the
nonlinear PDE. Appendix~\ref{app:coarse} supplies a broad sign proof
without the new direction quadrature or bump oracles.

For base lower endpoints $b_v^-,b_\Theta^-<0$ and slope upper endpoints
$l_v^+,l_\Theta^+<0$, choose the least nonnegative integer $n$ such that
\begin{equation}
 L=10^n>\max\{(-b_v^-)/(-l_v^+),
                         (-b_\Theta^-)/(-l_\Theta^+)\},\qquad\lambda=-L.
 \label{eq:rule}
\end{equation}
This algebraic rule was fixed before the direction evaluation. It alters
no direction, normalization, partition, target, or sign gate. The exact
threshold lies below $5.184\,10^{20}$ and above $10^{20}$, so it selects
$L=10^{21}$. Substitution in the exact affine relations
\eqref{eq:affine} proves the positive intervals in
Theorem~\ref{thm:main}. The coarse certificate independently selects the
same power and proves $v_{2,-10^{21}}>7\,10^{16}$ and
$\Theta_{2,-10^{21}}>5\,10^{21}$.

\section{Meaning, limits, and relation to prior work}

At this clock $V_0>0$ and
$\Theta_0=3\rho_b^{-1}\tanh(\ell T)>0$. All first integrated
coefficients vanish by the shared quadrupole mean. For each of the two
fixed preparations, $C^2$ parameter dependence gives
\begin{equation}
 \frac{V_\eps}{V_0}=1+\tfrac12\eps^2v_2+o(\eps^2),\qquad
 \Theta_\eps=\Theta_0+\tfrac12\eps^2\Theta_2+o(\eps^2).
 \label{eq:qualitative}
\end{equation}
For sufficiently small nonzero $\eps$ of either sign the base preparation
therefore has smaller positive volume and slower positive expansion than
the background. The second preparation has larger volume and faster
positive expansion. Both families remain expanding. A common qualitative
interval for these two fixed families is obtained by taking the smaller
of their unknown intervals. Opposite second coefficients do not establish
contraction, an effective physical amplitude, or a finite-time sign history.
Large second coefficients in this normalization are not finite-amplitude
field values, and no SI calibration is supplied.

Each nonzero slope permits arbitrary individual real second coefficients
as fixed finite $\lambda$ varies. This one line reverses both signs
because its slopes have the same sign. It does not prove rank two or
arbitrary independent pairs $(v_2,\Theta_2)$. Directions outside the
target causal region can be invisible there. The result also supplies no
full angular waveform, physical packet-generation process, hot matter,
radiation epoch, or evidence concerning a cosmological origin.

The general homogeneous freedom and affine second-response structure are
established machinery, not a new principle. Minamitsuji, Himemoto and
Sasaki explicitly retain initial-condition information in a second-order
scalar-induced Weyl term \cite{Minamitsuji}; their scalar does not couple
to the tension and their late-time regime differs. Deffayet's linear
master initial-boundary analysis \cite{Deffayet} explains why brane rows
do not replace bulk initial data. Hiramatsu and Koyama explicitly caution
that their curvature spectrum uses a special bulk preparation
\cite{Hiramatsu}. These works do not certify the compact direction or
the proper-volume sign here. A recent AdS$_3$/BCFT model illustrates
classical trajectory-data freedom \cite{Fujiki2025}; another proposes
Euclidean no-boundary regularity as a preparation prescription
\cite{Fujiki2026}. Both have a brane-local scalar and pure bulk gravity,
so their prescription cannot be imported into this bulk-scalar model
without a separate derivation.

The contribution is a conditional, model-specific quantitative
counterexample to simultaneous sign universality under the stated first
data. The focused primary-source comparison is not exhaustive and does
not establish originality or priority. The original result for its
expressly selected preparation remains valid. A predictive blast model
requires an additional physical rule selecting the second preparation;
this calculation neither supplies that rule nor disproves every blast
model. Other models require their own equations, initial conditions,
and observables.

\appendix
\section{Elementary sign certificate without direction quadrature}
\label{app:coarse}

\begin{lemma}[Coarse sufficient bounds]
For direction \eqref{eq:direction}, assume the retained coefficient
bounds and the following uniform exact caps on the causal initial segment:
\begin{gather}
 129\le\rho,\rho_b\le131,\quad
 3/25000\le-p\le1/7000,\quad 111/1000\le a\le14/125,\nonumber\\
 |p_y|\le1/4000,\quad |a_y|\le1/10000,\qquad
 1/1200\le-F'_b\le1/1000. \label{eq:coarsecaps}
\end{gather}
Then the full coupled slopes obey
\begin{equation}
 \mathcal L_M[h]\in[1.269\,10^9,2.198\,10^{10}],\quad
 \mathcal L_v[h]\in[-0.005883,-0.0001282],\quad
 \mathcal L_\Theta[h]\in[-359.5,-14.21]. \label{eq:coarseslopes}
\end{equation}
The same fixed $\lambda=-10^{21}$ makes both second coefficients positive.
\end{lemma}

\begin{proof}
For $t=r-1$, $D=t(1-t)$, on $1<r\le5/4$ direct differentiation gives
\begin{equation}
 Q'=\frac{Q(1-2t)}{D^2},\quad
 Q''=\frac{Q(1-6D+6D^2)}{D^4},\quad
 Q'''=\frac{Q(1-2t)(1-10D+24D^2-12D^3)}{D^6}.
 \label{eq:Qderivs}
\end{equation}
Here $0<D\le3/16$. The second numerator decreases from $1$ to
$11/128>0$; the cubic numerator decreases from $1$ to $-113/1024$.
Its absolute value is at most $1$. Set $w=1/D\ge16/3$.
For powers $2$ and $4$, $e^{-w}w^j$ decreases on this interval;
for power $6$ its maximum occurs at $w=6$. Rational Taylor bounds
$1/4<e^{-4/3}<4/15$ and $e^{-2}6^6<6400$ give, continuously through
the flat endpoint,
\begin{equation}
 0\le Q\le1/3,\quad |Q'|\le8,\quad |Q''|\le216,\quad
 |Q'''|\le6400. \label{eq:bumpcaps}
\end{equation}
No unresolved sampled tail enters these bounds. Moreover
$Q,Q',Q''$, and hence $s$, are positive for $1<r\le T$.
With $c_{\min}=22540234500000000/294499921$, the caps imply
$c\ge c_{\min}$. The exact endpoint values are
$Q(T)=e^{-4/3}$, $Q'(T)=(128/9)e^{-4/3}$, and
$Q''(T)=(5632/81)e^{-4/3}$. Positivity and integration by parts yield
\begin{equation}
 J_v>c_{\min}/4,\qquad
 J_\Theta>32c_{\min}/9,\qquad s(T)>1408c_{\min}/81. \label{eq:freegates}
\end{equation}
These avoid cancellation in the free moments.

For reproducible upper bounds put
\begin{gather*}
 R_+=(131/129)^3,\quad b_+=\ell\,131(14/125),\quad
 d_+=\ell\,131(1/4000)/(3/25000),\\
 c_+=28101137500000000/276922881,\\
 e_+=\ell^2\,131^2[(14/125)^2+1/10000],\qquad
 S_+=c_+[216+32b_++\ell^2],\\
 \begin{aligned}
 S_{r,+}=c_+\{&6400+864b_++24\ell^2\\
                  &+(4b_++d_+)[216+32b_++\ell^2]+32e_+\}.
 \end{aligned}
\end{gather*}
Equations \eqref{eq:profiles}, with the triangle inequalities
$\sqrt{4h^2+s^2}\le2|h|+|s|$ and $\sqrt2|V_L|\le2|V_L|$,
permit
\begin{equation*}
 D_0^{\rm init}\le16R_++S_+,\quad
 D_r^{\rm init}\le2R_+(216+24b_+)+S_{r,+},\quad
 D_v^{\rm init}\le2R_+(216+32b_+).
\end{equation*}
Substitution in the exact \eqref{eq:error} gives
$E_h\le243841044.434$. The strict positive gates for the coupled
histories are
\begin{equation}
 m(T)>1408c_{\min}/81-\Delta E_h>0,\quad
 \int_1^Tkm>(1-\eta)c_{\min}/4-\Delta^3E_h/6>0,
 \label{eq:coarsegates}
\end{equation}
and $\int_1^Tk_Tm>(1-\eta)32c_{\min}/9-\Delta^2E_h/2>0$.
The latter uses the exact upper bound $k_T\le1$ for the feedback error.
The matching upper bounds use \eqref{eq:bumpcaps} and the same errors.
All are rational operations apart from the stated rational exponential
bounds. Inserting \eqref{eq:coarsecaps} into \eqref{eq:slopes} yields
\eqref{eq:coarseslopes}. The exact coarse threshold in \eqref{eq:rule}
lies between $10^{20}$ and $6.345\,10^{20}$, selecting $10^{21}$ and
giving $v_{2,-10^{21}}>7\,10^{16}$ and
$\Theta_{2,-10^{21}}>5\,10^{21}$.
\end{proof}

The standalone coarse arithmetic reads only the two pinned retained
background/base-response JSON inputs. It calls no bump, background, or
evolved-first-field oracle and performs no direction quadrature. Its
agreement with the tighter calculation is an independent bounded method
under the same analytic and background premises, not a new PDE observation.
Source pins, the exact coarse report, the frozen whole-cell report,
decimal enclosure checks, and build receipt accompany this manuscript
\cite{records}.

\begin{thebibliography}{10}
\footnotesize
\setlength{\itemsep}{2pt}
\bibitem{Minamitsuji} M. Minamitsuji, Y. Himemoto, and M. Sasaki,
\emph{Geometry and cosmological perturbations in the bulk inflaton model},
Phys. Rev. D \textbf{68}, 024016 (2003),
\href{https://arxiv.org/abs/gr-qc/0303108}{arXiv:gr-qc/0303108}.
\bibitem{Deffayet} C. Deffayet,
\emph{A note on the wellposedness of scalar brane world cosmological perturbations},
Phys. Rev. D \textbf{71}, 023520 (2005),
\href{https://arxiv.org/abs/hep-th/0409302}{arXiv:hep-th/0409302}.
\bibitem{Hiramatsu} T. Hiramatsu and K. Koyama,
\emph{Numerical study of curvature perturbations in a brane-world inflation at high-energies},
JCAP \textbf{12} (2006) 009,
\href{https://arxiv.org/abs/hep-th/0607068}{arXiv:hep-th/0607068}.
\bibitem{Bruni} M. Bruni, S. Matarrese, S. Mollerach, and S. Sonego,
\emph{Perturbations of spacetime: gauge transformations and gauge invariance
at second order and beyond}, Class. Quantum Grav. \textbf{14} (1997),
\href{https://arxiv.org/abs/gr-qc/9609040}{arXiv:gr-qc/9609040}.
\bibitem{Mars} M. Mars, \emph{First and Second Order Perturbations of Hypersurfaces},
Class. Quantum Grav. \textbf{22} (2005),
\href{https://arxiv.org/abs/gr-qc/0507005}{arXiv:gr-qc/0507005}.
\bibitem{SMS} T. Shiromizu, K. Maeda, and M. Sasaki,
\emph{The Einstein Equations on the 3-Brane World}, Phys. Rev. D
\textbf{62}, 024012 (2000),
\href{https://arxiv.org/abs/gr-qc/9910076}{arXiv:gr-qc/9910076}.
\bibitem{Fujiki2025} K. Fujiki, H. Kanda, M. Kohara, and T. Takayanagi,
\emph{Brane Cosmology from AdS/BCFT} (2025),
\href{https://arxiv.org/abs/2501.05036}{arXiv:2501.05036}, Sec.~5.5.
\bibitem{Fujiki2026} K. Fujiki, M. Kohara, J. Kristiano, Y. Suzuki, and T. Takayanagi,
\emph{Effective Dynamics of Inflationary End-of-the-World Branes in AdS$_3$}
(2026), \href{https://arxiv.org/abs/2609.11643}{arXiv:2609.11643}.
\bibitem{records} R. Maldonado, \emph{HDBLAST conditional technical records
and reproducibility materials} (2026),
\href{https://github.com/maldonado-research/HDblast}{repository}.
Retained records include the localized initial-data and common-gauge
development, the 1 November prepared-shell volume capsule, and the
4 November preparation-difference capsule accompanying this working paper.
The latter contains the frozen experiment, exact rational outputs,
independent reconstructions, elementary sign certificate, and focused
primary-source literature comparison.
\end{thebibliography}
\end{document}
